\uuid{KbCU}
\exo7id{7894}
\titre{exo7 7894}
\auteur{mourougane}
\organisation{exo7}
\datecreate{2021-08-11}
\isIndication{false}
\isCorrection{false}
\chapitre{Groupe orthogonal et quaternions}
\sousChapitre{Groupe orthogonal et quaternions}
\cpge{mpsi-mp2i.calcul-matriciel-systemes.e, pcsi.determinants.c, ptsi.determinants.a, mpsi-mp2i.espaces-vectoriels.a, pcsi.espaces-vectoriels.a, ptsi.geometrie-espace.c, mpsi-mp2i.changements-bases.c, mpsi-mp2i.espaces-prehilbertiens.a, pcsi.espaces-prehilbertiens.a, ptsi.geometrie-plan.b, bcpst1.matrices, bcpst1.geometrie.f, pc.complements-espaces-vectoriels.c, psi.complements-espaces-vectoriels.c, pt.complements-algebre-lineaire.c, psi.espaces-prehilbertiens.a, pt.structure-prehilbertienne.a, bcpst2.espaces-vectoriels.a, bcpst2.produit-scalaire.a}
\module{Algèbre et géométrie}
\niveau{L3}
\difficulte{}

\contenu{
\begin{enumerate}
    \item \question{Déterminer la nature de l'espace vectoriel $V$ des matrices anti-hermitiennes (i. e. $M^\star:=-M$) de $M(2,\Cc)$ de trace nulle.}
    \item \question{Écrire la forme générale d'une matrice de $V$ à l'aide de trois nombres réels. En déduire une base de $V$.}
    \item \question{Montrer que $\ll P,P'\gg:=-\frac{1}{2}trace(PP')$ définit un produit scalaire sur $V$.}
\end{enumerate}
}
